% RSE abstract (at most 250 words; no citations; headline first; review S3, P10, C6).
% First revision (RSASE-D-26-03629, Reviewer 1 comment 1): the data analysed and the limitations
% are now stated; the flooded-area bias and the Great Salt Lake distances are in the Results only.
\begin{abstract}
Projections of the seabed that shrinking seas expose rely on global bathymetry grids that are rarely tested against seabed already emerged. \rev{We tested five grids (GEBCO\_2026, GEBCO\_2022, ETOPO~2022, ETOPO1 and SRTM15+) in the North Caspian Sea against shorelines observed from 1992 to 2026. We classified \metric{erl.text.inventory.n_landsat} Landsat scenes with MNDWI and Otsu thresholds, levelled each waterline with DAHITI altimetry, measured a noise floor from two halves of the imagery, and checked the result against ICESat-2 heights.} Flooded from the open sea at each observed level, four grids misplace the shoreline in \metric{erl.text.h1.majority.n_sectors_min} or \metric{erl.text.h1.majority.n_sectors_max} of seven sectors, even after \rev{an area-fitted vertical offset}. On the Kizlyar--Dagestan and Volga--Ural interfluve coasts, where ICESat-2 confirms the observations, they miss by \metricdp{rev.s.headline.west.d_km_min}{1}--\metricdp{rev.s.headline.west.d_km_max}{1}~km, \metric{rev.s.headline.west.dq_median_min}--\metric{rev.s.headline.west.dq_median_max} times the noise floor, or \metricdp{rev2.d.vert.westinterfluve.four_Dg_min_m}{1}--\metricdp{rev2.d.vert.westinterfluve.four_Dg_max_m}{1}~m in height. Whole-metre storage, coarse cells, level errors and datum choices do not explain \rev{the misses,} no single vertical shift corrects them\rev{, and they persist without the river-influenced Volga and Ural delta sectors}. ETOPO~2022 passes only where \rev{it holds} land heights measured after emergence. \rev{A waterline elevation model was no more accurate than the grids. The conclusions rest on three coasts of one basin: elsewhere grid error cannot be separated from observation error. Under reeds the Landsat shoreline lies \metric{rev3.lidar.both.checked.vegetated.median_h_abs_km.min}--\metric{rev3.lidar.both.checked.vegetated.median_h_abs_km.max}~km seaward of the lidar ground, and wind setup widens the noise floor. Areas the grids place in the first metre of further fall (\metricsf{rev2.n.reref.area.main.E.1.domain.min_km2}{2}--\metricsf{rev2.n.reref.area.main.E.1.domain.max_km2}{3}~km\textsuperscript{2}) are therefore conditional, not predictions.}
\end{abstract}
